%!TEX root = ../chapter04.tex 
%----------------------------------------------------------------------------------------
%	
%----------------------------------------------------------------------------------------


%----------------------------------------------------------------------------------------
%	 
%----------------------------------------------------------------------------------------
\section{Facteurs, multiples et division euclidienne}
\begin{definition}[Multiples]
	Soit $n$ et $m\in\mathbb{N}$.
	
	$m$ est un multiple de $n$ \textbf{si et seulement si} il existe $k\in\mathbb{N}$ tel que $m=k\times n$.
	
	Les multiples de $n$ tous les nombres de la forme $kn$, où $k\in\mathbb{N}$.
\end{definition}
\begin{example}\hfill 
	\begin{enumerate}[label={\alph*)}, leftmargin=*]  
		\item $45$ est un multiple de $9$ et de $5$ car $45=5\times 9$.  
		\item Enumérer tous les diviseurs de $45$. 
	\end{enumerate}
\end{example}
\begin{remark}
	Pour tout $m\in \mathbb{N}$,  $0=0\times m$. Donc  \textbf{$0$ est un multiple de tout entier, et tout entier divise $0$}. Néanmoins, l'usage est souvent d'ignorer $0$ comme multiple. 
\end{remark}
\begin{definition}[Division entière] Soit $n\in\mathbb{N}$ et $d\in\mathbb{N}^*$.
	
	On appelle quotient de $n$ par $d$ l'entier $q$ tel que : 
	\[ qd\leqslant n < qd + d\]
	
	Le reste de la division de $n$ par $d$  est l'entier définit par :
	\[ n = qd + r
	\qquad 0\leqslant r < d \]
	
	$n$ est un multiple de $d$ s.s.i. le reste de la division de $n$
	par $d$ est nul.
\end{definition}
\begin{example} 
	\begin{marginfigure} 
		\opidiv[displayintermediary=all,voperation=top]{17}{5}\hfill 
		\opidiv[displayintermediary=all,voperation=top]{20}{5}
		\caption{$17=3\times 5 + 2$ et $20 = 4 \times 5$}
		\label{chapitre03:figure:divisionentiere}
	\end{marginfigure}
	\begin{enumerate}[label={\alph*)}, leftmargin=*]  
		\item $3\times 5\leqslant 17<4\times 5$ et $17=3\times 5+2$. 
		\item $4\times 5\leqslant 20<5\times 5$ et $20=4\times 5+0$.  
	\end{enumerate} 
\end{example}
%----------------------------------------------------------------------------------------
%	EXERCICES
%----------------------------------------------------------------------------------------

\newpage 
\loadgeometry{widelayout}  
\pagestyle{scrheadings} % Print headers again  
\KOMAoptions{
	headwidth=textwithmarginpar,
	headsepline=true,
	headsepline= 0.5pt,
	footwidth=textwithmarginpar,
} 
\onehalfspacing


\subsection{Exercices}
\begin{exercise}\hfill 
	\begin{enumerate}[label={\arabic*)}, leftmargin=*]  
		\item  Pour $n\in\mathbb{N}$, on pose $A(n)=n^2+(n+1)^2$.
		\begin{enumerate}[label={\alph*)}, leftmargin=*]  
			\item Vérfiez que si $n=5$, alors $A(n)$ est impair.
			\item Montrer que pour tout $n\in\mathbb{N}$, $A(n)$ est impair
		\end{enumerate} 
		\item Montrer que pour tout $n\in\mathbb{N}$ , $B(n)=(3n+1)^2-(3n-1)^2$ est un multiple de $4$.
		\item Montrer que pour tout $n\in\mathbb{N}$ , $C(n)=(2n+3)^2-(2n-3)^2$ est un multiple de $8$.
	\end{enumerate}
\end{exercise}
\begin{example} Pour $n\in\mathbb{N}$. \\
	Exprimer à l'aide de $n$, le quotient et le reste de la division de $a$ par $b$ dans les cas suivants :
	\begin{multicols}{3}
		\begin{enumerate}[label={\arabic*)}, leftmargin=*]  
			\item$a=15n+2$ et $b=3$.
			\item$a=15n+13$ et $b=5$.
			\item$a=15n-4$ et $b=5$.
		\end{enumerate} 
	\end{multicols}
	\vspace{2cm} 
\end{example}
\begin{exercise}[À vous] Pour $n\in\mathbb{N}$. Exprimer à l'aide de $n$, le quotient et le reste de la division de $a$ par $b$ dans les cas suivants :
	\begin{multicols}{3} 
		\begin{enumerate}[label={\arabic*)}, leftmargin=*]  
			\item $a=6n+2$ et $b=2$
			\item $a=12n+2$ et $b=3$
			\item $a=12n+7$ et $b=3$ 
			\item $a=6n-5$ et $b=6$ 
			\item $a=12n-7$ et $b=6$ 
			\item $a=6n-5$ et $b=3$  
		\end{enumerate}  
	\end{multicols}
\end{exercise} 
\begin{exercise}  Mêmes consignes.
	\begin{multicols}{3} 
		\begin{enumerate}[label={\arabic*)}, leftmargin=*]   
			\item$a=(2n+1)^2$ et $b=4$
			\item$a=(2n-1)^2$ et $b=2$
			%			\item $a=(n+1)^2-n^2$ et $b=2$ 
			\item $a=(3n+2)^2$ et $b=3$ 
		\end{enumerate}  
	\end{multicols}
\end{exercise}
\begin{exercise}
	Associe pour les énoncés le mieux adapté.
	
	\noindent\Relie[Ecart=1cm,LargeurG=0.43\linewidth,LargeurD=0.5\linewidth]{%Solution,
		{Pour tout $p, q\in\mathbb{N}$, $2p+1$ et $2p+3$ sont \dots}/deux nombres consécutifs/3,
		{Pour tout $p, q\in\mathbb{N}$, $2p+1$ et $2p+2$ sont \dots}/deux nombres pairs consécutifs/5,
		{Pour tout $p, q\in\mathbb{N}$, $2p$ et $2q$ sont \dots}/deux nombres impairs consécutifs/7,
		{Pour tout $p, q\in\mathbb{N}$, $2p$ et $2p+2$ sont \dots}/un nombre pair et le nombre impair suivant/2,
		{Pour tout $p, q\in\mathbb{N}$, $p$ et $p+1$ sont \dots}/un nombre impair et le nombre pair suivant/1,
		{Pour tout $p, q\in\mathbb{N}$, $2p$ et $2p+1$ sont \dots}/deux nombres impairs quelconques/4,
		{Pour tout $p, q\in\mathbb{N}$, $2p+1$ et $2q+1$ sont \dots}/deux nombres pairs quelconques/6}
\end{exercise}
\begin{exercise}\hfill 
	\begin{enumerate}[label={\arabic*)}, leftmargin=*]   
		\item Montrer que la somme de 3 nombres consecutifs est toujours un multiple de 3.
		\item Montrer que la somme de 3 nombres pairs consécutifs est toujours un multiple de 6.
		\item Montrer que la somme de 4 nombres impairs consécutifs est toujours un multiple de 8.
	\end{enumerate} 
\end{exercise}
%----------------------------------------------------------------------------------------
%	
%----------------------------------------------------------------------------------------